\documentclass[12pt,a4paper]{article}
%
%	STARWARE - Stile base per la scrittura di documenti interni in LaTex
%

%
%	Pacchetti globali
%	Fare qui eventuali aggiunte!
%
\usepackage[utf8]{inputenc}
\usepackage[italian]{babel}
\usepackage[babel]{csquotes}
\usepackage{url}
\usepackage{graphicx}
\usepackage[colorlinks]{hyperref}
\usepackage{lastpage}
\usepackage{fancyhdr}
\usepackage[top=1cm,bottom=4cm,left=80pt,right=80pt]{geometry} %disegna la linea
\usepackage{listings} %per grandi porzioni di codice
\usepackage{color}
\usepackage[table]{xcolor}
\usepackage{booktabs,tabularx}
\usepackage{makeidx}
\usepackage{fixltx2e}
\usepackage{hyperref}
\usepackage{enumitem}
\usepackage{color}
\usepackage[T1]{fontenc}
\usepackage{float}
\usepackage{svg}
\usepackage{amsmath}
\usepackage[toc]{glossaries}
\usepackage{dirtree}
\usepackage{listings}
\usepackage{eurosym}

\makeglossaries

\bibliographystyle{alpha}

%
%	VARIABILI GLOBALI
%
\newcommand{\nomeGruppo}{StarWare}
\newcommand{\mailGruppo}{starware.swe@gmail.com}
\newcommand{\uni}{Universit\`{a} degli Studi di Padova}
\newcommand{\uniAA}{2015/2016}
\newcommand{\Cardin}{Prof. Riccardo Cardin}
\newcommand{\Vardanega}{Prof. Tullio Vardanega}
\newcommand{\Zucchetti}{Zucchetti S.p.a.}
\newcommand{\prj}{Quizzipedia}
\newcommand{\prjL}{Quizzipedia: software per la gestione di questionari}

\newcommand{\AVI}{Alessio Vitella}
\newcommand{\AVE}{Andrea Venier}
\newcommand{\NDC}{Nicola De Cao}
\newcommand{\IB}{Igor Baylyak}
\newcommand{\WS}{Walter Sandon}
\newcommand{\TP}{Thomas Pigarelli}
\newcommand{\AB}{Anna Bonaldo}

\newcommand{\mgls}[1]{\gls{#1}\textsubscript{G}}
\newcommand{\mglspl}[1]{\glspl{#1}\textsubscript{G}}
\newcommand{\mGls}[1]{\Gls{#1}\textsubscript{G}}
\newcommand{\mGlspl}[1]{\Glspl{#1}\textsubscript{G}}

\newcommand{\AM}{\emph{\mGls{amministratore}}}
\newcommand{\AN}{\emph{\mGls{analista}}}
\newcommand{\PG}{\emph{\mGls{progettista}}}
\newcommand{\PR}{\emph{\mGls{programmatore}}}
\newcommand{\VR}{\emph{\mGls{verificatore}}}
\newcommand{\RE}{\emph{\mGls{responsabile}}}
\newcommand{\PM}{\RE}

\newcommand{\AMpl}{\emph{\mGlspl{amministratore}}}
\newcommand{\ANpl}{\emph{\mGlspl{analista}}}
\newcommand{\PGpl}{\emph{\mGlspl{progettista}}}
\newcommand{\PRpl}{\emph{\mGlspl{programmatore}}}
\newcommand{\VRpl}{\emph{\mGlspl{verificatore}}}
\newcommand{\REpl}{\emph{\mGlspl{responsabile}}}
\newcommand{\PMpl}{\REpl}

\newcommand{\RR}{\emph{\mGls{revisione dei requisiti}}}
\newcommand{\RA}{\emph{\mGls{revisione di accettazione}}}
\newcommand{\RP}{\emph{\mGls{revisione di progettazione}}}
\newcommand{\RQ}{\emph{\mGls{revisione di qualifica}}}

\newcommand{\NdP}{\emph{\mGls{norme di progetto}}}
\newcommand{\SdF}{\emph{\mGls{studio di fattibilita}}}
\newcommand{\AdR}{\emph{\mGls{analisi dei requisiti}}}
\newcommand{\PdP}{\emph{\mGls{piano di progetto}}}
\newcommand{\PdQ}{\emph{\mGls{piano di qualifica}}}

\newcommand{\latex}[1]{\texttt{#1}}
\newcommand{\fileName}[1]{\texttt{#1}}
\newcommand{\filePath}[1]{\texttt{#1}}
\newcommand{\TODO}[1]{\texttt{\large \color{red} \underline{TODO: #1}}}

\newcommand{\licenza}{GNU GENERAL PUBLIC LICENSE V2}

%per compilare il template usare questi sotto e commentare glia altri:
%\newcommand{\logoLungo}{../imgs/logoLungo.png}
%\newcommand{\logoGrande}{../imgs/logoGrande.png}
\newcommand{\logoLungo}{../../../template/imgs/logoLungo.png}
\newcommand{\logoGrande}{../../../template/imgs/logoGrande.png}

%
%	Setup stili
%

\newcommand{\HRule}{\rule{\linewidth}{0.5mm}}

\definecolor{dkgreen}{rgb}{0,0.6,0}
\definecolor{gray}{rgb}{0.5,0.5,0.5}
\definecolor{mauve}{rgb}{0.58,0,0.82}
\definecolor{light}{RGB}{255,255,190}

%
%	Setup di pagina
%

%colorazione link
\hypersetup
{
	colorlinks=true,
	linkcolor=black,
	urlcolor=blue,
	citecolor=blue
}

%	Setup Header + Footer

\pagestyle{fancy}
\setlength{\headheight}{2cm} %settato grandezza header

\renewcommand{\footrulewidth}{0.5pt} %ridefinisco il valore della riga di intestazione
\renewcommand{\headrulewidth}{0.5pt} %ridefinisco il valore della riga di pie' di pagina
\addtolength{\headwidth}{\marginparsep}
\addtolength{\headwidth}{\marginparwidth}

\fancyhead{} %annulla head di default
\fancyfoot{} %annulla foot di default

%	Logo intestazione
\lhead{\includegraphics[scale=0.12]{\logoLungo}}


%	footer
\cfoot{
	\uni \ - \uniAA \\
	\href{mailto:\mailGruppo}{\mailGruppo}\\
	{\tiny Questo documento è distribuito sotto licenza {\licenza}}
}
\rfoot{
	\thepage\ di \pageref{LastPage}
}


% test subsubsubsection

\usepackage{titlesec}
\usepackage{hyperref}

\titleclass{\subsubsubsection}{straight}[\subsection]

\newcounter{subsubsubsection}[subsubsection]
\renewcommand\thesubsubsubsection{\thesubsubsection.\arabic{subsubsubsection}}
\renewcommand\theparagraph{\thesubsubsubsection.\arabic{paragraph}} % optional; useful if paragraphs are to be numbered

\titleformat{\subsubsubsection}
  {\normalfont\normalsize\bfseries}{\thesubsubsubsection}{1em}{}
\titlespacing*{\subsubsubsection}
{0pt}{3.25ex plus 1ex minus .2ex}{1.5ex plus .2ex}

\makeatletter
\renewcommand\paragraph{\@startsection{paragraph}{5}{\z@}%
  {3.25ex \@plus1ex \@minus.2ex}%
  {-1em}%
  {\normalfont\normalsize\bfseries}}
\renewcommand\subparagraph{\@startsection{subparagraph}{6}{\parindent}%
  {3.25ex \@plus1ex \@minus .2ex}%
  {-1em}%
  {\normalfont\normalsize\bfseries}}
\def\toclevel@subsubsubsection{4}
\def\toclevel@paragraph{5}
\def\toclevel@paragraph{6}
\def\l@subsubsubsection{\@dottedtocline{4}{7em}{4em}}
\def\l@paragraph{\@dottedtocline{5}{10em}{5em}}
\def\l@subparagraph{\@dottedtocline{6}{14em}{6em}}
\makeatother

\setcounter{secnumdepth}{4}
\setcounter{tocdepth}{4}


%pdflatex -synctex=1 -interaction=nonstopmode %.tex|makeglossaries %|pdflatex -synctex=1 -interaction=nonstopmode %.tex|pdflatex -synctex=1 -interaction=nonstopmode %.tex
\makeglossary

\newglossaryentry{responsabile} {
	name=Responsabile,
	description={è il responsabile della gestione, pianificazione e realizzazione del progetto.},
	plural=Responsabili
}

\newglossaryentry{verificatore} {
	name=Verificatore,
	description={è il responsabile dell'attività di verifica.},
	plural=Verificatori
}

\newglossaryentry{programmatore} {
	name=Programmatore,
	description={è responsabile delle attività di codifica miranti alla realizzazione del prodotto e delle componenti di ausilio necessarie per l'esecuzione delle prove di verifica e validazione.},
	plural=programmatori
}

\newglossaryentry{progettista} {
	name=Progettista,
	description={è responsabile delle attività di progettazione.},
	plural=Progettisti
}

\newglossaryentry{analista} {
	name=Analista,
	description={è responsabile delle attività di analisi. },
	plural=Analisti
}

\newglossaryentry{amministratore} {
	name=Amministratore,
	description={è responsabile dell'efficienza e dell'operatività dell'ambiente di sviluppo; si occupa della redazione e attuazione di piani e procedure di gestione della qualità; inoltre gestisce l'archivio della documentazione del progetto.},
	plural=Amministratori
}

\newglossaryentry{revisione dei requisiti} {
	name=Revisione dei requisiti,
	description={è una revisione formale che determina l'accesso del gruppo al progetto didattico e la concordanza con il cliente di una visione condivisa del prodotto atteso.}
}

\newglossaryentry{revisione di accettazione} {
	name=Revisione di accettazione,
	description={è una revisione formale per l'accertamento del soddisfacimento di tutti i requisiti e il completamento del progetto.}
}

\newglossaryentry{revisione di progettazione} {
	name=Revisione di progettazione,
	description={è una revisione di progresso che accerta la realizzabilità del prodotto e informa il cliente sulle caratteristiche del prodotto.}
}

\newglossaryentry{revisione di qualifica} {
	name=Revisione di qualifica,
	description={è una revisione di progresso che approva l'esito finale delle verifiche e attiva la fase di validazione.}
}

\newglossaryentry{analisi} {
	name=Analisi,
	description={è il periodo di preparazione e produzioni di documenti che precede la Revisione dei requisiti.}
}

\newglossaryentry{progettazione architetturale e di dettaglio} {
	name=Progettazione architetturale e di dettaglio,
	description={è il periodo che intercorre tra la Revisione dei requisiti e la Revisione di progettazione.}
}

\newglossaryentry{codifica} {
	name=Codifica,
	description={è il periodo che intercorre tra la Revisione di progettazione e la Revisione di qualifica.}
}

\newglossaryentry{verifica e validazione} {
	name=Verifica e validazione,
	description={è il periodo che intercorre tra la Revisione di qualifica e la Revisione di accettazione.}
}

\newglossaryentry{repository} {
	name=Repository,
	description={è dove i file sono memorizzati, spesso su un server.}
}

\newglossaryentry{ruolo} {
	name=Ruolo,
	description={una delle  gure professionali che una persona  sica interpreta nel corso del progetto. I ruoli sono: responsabile, amministratore, analista, progettista, programmatore e veri catore.},
    plural=ruoli
}

\newglossaryentry{SVG} {
	name=SVG,
	description={è un formato per la visualizzazione di oggetti in grafica vettoriale. Per maggiori informazioni si veda \href{https://it.wikipedia.org/wiki/Scalable_Vector_Graphics}{qui}.}
}

\newglossaryentry{PNG} {
	name=PNG,
	description={abbreviazione di Portable Network Graphics, è un formato di file per memorizzare immagini. Per ulteriori informazioni si veda http://it.wikipedia.org/wiki/Portable_Network_Graphics.}
}

\newglossaryentry{PDF} {
	name=PDF,
	description={è un formato di file basato su un linguaggio di descrizione di pagina sviluppato da Adobe Systems nel 1993 per rappresentare documenti in modo indipendente dall’hardware e dal software utilizzati per generarli o per visualizzarli. Per ulteriori informazioni si veda http: //it.wikipedia.org/wiki/Portable_Document_Format.}
}

\newglossaryentry{UML} {
	name=UML,
	description={è un linguaggio di modellazione e specifica basato sul paradigma object-oriented. Per ulteriori informazioni si veda http://it.wikipedia.org/wiki/Unified_Modeling_Language.}
}

\newglossaryentry{walkthrough} {
	name=Walkthrough,
    description={consiste nella lettura di un documento o codice cercando errori ed anomalie senza un'idea precisa di quali tipi di errori sarà possibile trovare.}
}

\newglossaryentry{lista di controllo} {
	name=Lista di controllo,
	description={è un elenco di cose da fare per eseguire una determinata attività.}
}

\newglossaryentry{inspection} {
	name=Inspection,
	description={è la lettura mirata di un documento o codice cercando errori specifici.}
}

\newglossaryentry{milestone} {
	name=Milestone,
	description={momento saliente nello sviluppo di un prodotto software per la quale devono essere pronti documenti e/o funzionalità.}
}

\newglossaryentry{ticket} {
	name=Ticket,
	description={rappresenta un compito nell'organizzazione e distribuzione del lavoro all'interno del progetto.},
	plural=tickets
}

\newglossaryentry{commit} {
	name=Commit,
	description={è la copia di modifiche fatte su file locali verso la repository remota. Esso rappresenta anche un particolare stato della repository nel tempo.}
}

\newglossaryentry{versionamento} {
	name=Versionamento,
	description={è la gestione di un versioni multiple di un insieme di informazioni. Per maggiori informazioni si veda \href{http://it. wikipedia.org/wiki/Controllo_versione}{qui}.}
}

\newglossaryentry{task} {
	name=Task,
	description={è un compito secondo la definizione dello standard IEEE 12207.},
	plural=tasks
}

\newglossaryentry{attivita} {
	name=attivita,
	description={è un insieme di task.}
}

\newglossaryentry{redattore} {
	name=Redattore,
	description={colui che redige un documento.},
	plural=redattori
}

\newglossaryentry{proponente} {
	name=Proponente,
	description={colui che ha proposto al committente un capitolato d'appalto.}
}

\newglossaryentry{committente} {
	name=Committente,
	description={colui che assegna un compito. In questo caso è il Professor Tullio Vardanega.}
}

\newglossaryentry{quality assurance} {
	name=Quality assurance,
	description={è l'insieme delle attività volte a garantire il soddisfacimento degli obiettivi della qualità.}
}

\newglossaryentry{telegram} {
	name=Telegram,
	description={è un servizio di messaggistica istantanea utilizzato dal gruppo per comunicazioni interne. Per maggiori informazioni si veda \href{https://it.wikipedia.org/wiki/Telegram_(software)}{qui}.}
}

\newglossaryentry{browser} {
	name=Browser,
	description={è un'applicazione per il recupero, la presentazione e la navigazione di risorse web.}
}

\newglossaryentry{google drive} {
	name=Google Drive,
	description={è un servizio di memorizzazione e sincronizzazione online introdotto da Google il 24 aprile 2012. Per maggiori informazioni si veda \href{https://it.wikipedia.org/wiki/Google_Drive}{qui}.}
}

\newglossaryentry{skype} {
	name=Skype,
	description={è un software proprietario freeware di messaggistica istantanea e VoIP. Per maggiori informazioni si veda \href{https://it.wikipedia.org/wiki/Skype}{qui}.}
}

\newglossaryentry{gantt} {
	name=Gantt,
	description={è un diagramma di supporto alla gestione dei progetti.}
}

\newglossaryentry{projectlibre} {
	name=ProjectLibre,
	description={è un software di gestione progettuale.}
}

\newglossaryentry{PERT} {
	name=PERT,
	description={è uno strumento volto alla programmazione delle attività che compongono il progetto e, più in generale, alla gestione degli aspetti temporali di quest'ultimo.}
}

\newglossaryentry{subtask} {
	name=Subtask,
	description={è un task compreso all'interno di un altro task. La totalità di tutti i subtasks costituisce un intero task.}
	plural=subtasks
}

\newglossaryentry{ticketing} {
	name=Ticketing,
	description={procedura con la quale il Responsabile assegna un task.}
}

\newglossaryentry{git} {
	name=Git,
	description={è un sistema software di controllo di versione distribuito}
}

\newglossaryentry{Quizzpedia} {
	name=Quizzpedia,
	description={è il nome del prodotto software richiesto dal capitolato d'appalto scelto.}
}

\newglossaryentry{schierabile} {
	name=Schierabile,
	description={è la capacità di rilasciare al cliente, con relativa installazione e messa in funzione o esercizio, di una applicazione o di un sistema software tipicamente all'interno di un sistema informatico aziendale.}
}

\newglossaryentry{cross-platform} {
	name=Cross-platform,
	description={può essere riferito ad un linguaggio di programmazione, ad un'applicazione software o ad un dispositivo hardware che funziona su più di un sistema.}
}

\newglossaryentry{QML} {
	name=QML,
	description={è un "Domain Specific Language" richiesto dal capitolato d'appalto per la definizione delle domande all'interno del sistema.}
}

\newglossaryentry{Tomcat} {
	name=Tomcat,
	description={è un application server nella forma di contenitore servlet open source sviluppato dalla Apache Software Foundation. Per maggiori informazioni si veda \href{https://it.wikipedia.org/wiki/Apache_Tomcat}{qui}.}
}

\newglossaryentry{Java} {
	name=Java,
	description={è un linguaggio di programmazione orientato agli oggetti, specificatamente progettato per essere il più possibile indipendente dalla piattaforma di esecuzione. Per maggiori informazioni si veda \href{https://it.wikipedia.org/wiki/Java_(linguaggio_di_programmazione)}{qui}.}
}

\newglossaryentry{Node.js} {
	name=Node.js,
	description={\TODO{}}
}

\newglossaryentry{JavaScript} {
	name=JavaScript,
	description={\TODO{}}
}

\newglossaryentry{PostgreSQL} {
	name=PostgreSQL,
	description={\TODO{}}
}

\newglossaryentry{MongoDB} {
	name=MongoDB,
	description={\TODO{}}
}

\newglossaryentry{HTML5} {
	name=HTML5,
	description={\TODO{}}
}

\newglossaryentry{CSS3} {
	name=CSS3,
	description={\TODO{}}
}

\newglossaryentry{XML} {
	name=XML,
	description={\TODO{}}
}

\newglossaryentry{NoSQL} {
	name=NoSQL,
	description={\TODO{}}
}

\newglossaryentry{Scala} {
	name=Scala,
	description={\TODO{}}
}

\newglossaryentry{Akka} {
	name=Akka,
	description={\TODO{}}
}

\newglossaryentry{BLE} {
	name=BLE,
	description={\TODO{}}
}

\newglossaryentry{MQTT} {
	name=MQTT,
	description={\TODO{}}
}

\newglossaryentry{AWS} {
	name=AWS,
	description={\TODO{}}
}

\newglossaryentry{Heroku} {
	name=Heroku,
	description={\TODO{}}
}

\newglossaryentry{Github} {
	name=Github,
	description={\TODO{}}
}

\newglossaryentry{template} {
	name=template,
	description={\TODO{}}
	plural=templates
}

\newglossaryentry{Dropbox} {
	name=Dropbox,
	description={\TODO{}}
}

\newglossaryentry{open-source} {
	name=open-source,
	description={\TODO{}}
}

\newglossaryentry{project management} {
	name=project management,
	description={\TODO{}}
}

\newglossaryentry{Linux} {
	name=Linux,
	description={\TODO{}}
}

\newglossaryentry{Windows} {
	name=Windows,
	description={\TODO{}}
}

\newglossaryentry{Mac OS} {
	name=Mac OS,
	description={\TODO{}}
}

\newglossaryentry{program evaluation and review techniques} {
	name=program evaluation and review techniques,
	description={\TODO{}}
}

\newglossaryentry{schedule variance} {
	name=schedule variance,
	description={\TODO{}}
}

\newglossaryentry{cost variance} {
	name=cost variance,
	description={\TODO{}}
}


\newglossaryentry{merge} {
	name=merge,
	description={\TODO{}}
}

\newglossaryentry{slack} {
	name=slack,
	description={\TODO{}}
}

\newglossaryentry{baseline} {
	name=baseline,
	description={\TODO{}}
}

\newglossaryentry{Asana} {
	name=Asana,
	description={\TODO{}}
}

\newglossaryentry{deadline} {
	name=deadline,
	description={\TODO{}}
}

\newglossaryentry{revert} {
	name=revert,
	description={\TODO{}}
}

\newglossaryentry{backup} {
	name=backup,
	description={\TODO{}}
}

\newglossaryentry{push} {
	name=push,
	description={\TODO{}}
}

\newglossaryentry{teamwork} {
	name=teamwork,
	description={\TODO{}}
}

\newglossaryentry{evento} {
	name=evento,
	description={\TODO{}}
}

\newglossaryentry{etichetta} {
	name=etichetta,
	description={\TODO{}}
	plural=etichette
}

\newglossaryentry{bug} {
	name=bug,
	description={\TODO{}}
	plural=bugs
}

\newglossaryentry{software} {
	name=software,
	description={\TODO{}}
}

\newglossaryentry{desktop} {
	name=desktop,
	description={\TODO{}}
}

\newglossaryentry{draw.io} {
	name=draw.io,
	description={\TODO{}}
}

\newglossaryentry{definizione di prodotto} {
	name=definizione di prodotto,
	description={\TODO{}}
}

\newglossaryentry{specifica tecnica} {
	name=specifica tecnica,
	description={\TODO{}}
}

\newglossaryentry{manuale utente} {
	name=manuale utente,
	description={\TODO{}}
}

\newglossaryentry{tracy} {
	name=tracy,
	description={\TODO{}}
}

\newglossaryentry{package} {
	name=package,
	description={\TODO{}}
}

\newglossaryentry{stakeholder} {
	name=stakeholder,
	description={\TODO{}}
}

\newglossaryentry{branch} {
	name=branch,
	description={\TODO{}}
	plural=branches
}

\newglossaryentry{pull} {
	name=pull,
	description={\TODO{}}
}

\newglossaryentry{underscore} {
	name=underscore,
	description={\TODO{}}
}

\newglossaryentry{spelling} {
	name=spelling,
	description={\TODO{}}
}

\newglossaryentry{cloud} {
	name=cloud,
	description={\TODO{}}
}

\newglossaryentry{smartphone} {
	name=smartphone,
	description={\TODO{}}
	plural=smartphones
}

\newglossaryentry{checkbox} {
	name=checkbox,
	description={\TODO{}}
}

\newglossaryentry{appointed} {
	name=appointed,
	description={\TODO{}}
}

\newglossaryentry{makefile} {
	name=makefile,
	description={\TODO{}}
}

\newglossaryentry{gulpease} {
	name=Gulpease,
	description={\TODO{}}
}



%Titolo documento
\newcommand{\titoloDocumento}{Piano di Progetto}

%Prima data di creazione del documento
\newcommand{\dataCreazione}{24 dicembre 2015}

%Inserite la versione attuale del documento
\newcommand{\versione}{1.0.1}

%Stato in cui si trova il documento: Formale solo all'atto di consegna
\newcommand{\stato}{Informale}

%Uso del documento
\newcommand{\uso}{Esterno}

\rhead{\titoloDocumento}
\lfoot{Versione: \versione}
\title{\titoloDocumento}

\begin{document}
\begin{titlepage}
\begin{center}
\today \\
\vspace{1cm}
\begin{Huge}
\textbf{\nomeGruppo} \\
\end{Huge}
\textbf{\prjL} \\
\vspace{1cm}
\includegraphics[scale=0.5]{\logoGrande}
\vspace{1cm}

\HRule \\[0.4cm]
\begin{Huge}
{\huge \bfseries \titoloDocumento}\\[0.4cm]
\end{Huge}
\HRule \\[1cm]
\vfill

\begin{table}[h]
\begin{center}
\begin{tabular}{r | l}
\multicolumn{2}{c}{\textbf{Informazioni sul documento}}\\
\midrule
\textbf{Nome Documento}	&	\titoloDocumento	\\
\textbf{Versione}	&	\versione	\\
\textbf{Stato}	&	\emph{\stato}	\\
\textbf{Uso}	&	\emph{\uso}	\\
\textbf{Data Creazione}	&	\dataCreazione	\\
\textbf{Data Ultima Modifica}	&	\today	\\
\textbf{Redazione}	& \NDC \\
\ 	& \TP \\
\ 	& \IB \\
\textbf{Verifica}	&	Mr Q	\\
\textbf{Approvazione}	& \NDC \\
\textbf{Lista distribuzione}	&	\nomeGruppo	\\
\ 	&	\Vardanega	\\
\ 	&	\Cardin	\\
\ 	&	Il proponente \Zucchetti	\\

\end{tabular}
\end{center}
\end{table}

\end{center}
\end{titlepage}
\newpage

\Large{\textbf{Registro delle modifiche}}\\
\normalsize

\begin{table}[h]
\begin{center}

\begin{tabular}{p{0.12\textwidth} p{0.2\textwidth} p{0.18\textwidth} p{0.39\textwidth}}
\toprule
\textbf{Versione}	&	\textbf{Autore}	&	\textbf{Data}	&	\textbf{Descrizione}\\
\midrule
\midrule
1.0.1 & \NDC & 2015-12-27 &  Sistemato template e attività di Analisi \\
\midrule
1.0.0 & \NDC & 2015-12-24 &  Creazione documento \\
\bottomrule
\end{tabular}
\caption{Versionamento del documento}
\label{tabVers1}
\end{center}
\end{table}
\newpage

\tableofcontents
\newpage

\listoftables
\listoffigures
\newpage

\section{Introduzione}
\subsection{Scopo}
\subsection{Descrizione}

Si divide lo sviluppo del progetto in 5 fasi:
\begin{itemize}
	\item \textbf{Analisi:} (quella corrente) [Per RRequisiti]: studio di fattibilità, norme di progetto, analisi dei requisiti, piano di progetto, piano di qualifica, glossario (eventuale?), lettera di presentazione.
	\item \textbf{Analisi dettaglio:} integrazione e consolidazione dei requisiti. Progettazione architetturale [Per RProgettazione]: Viene redatta la specifica tecnica che descrive le scelte architetturali per lo sviluppo del prodotto (design pattern, diagrammi, etc...); tutti gli altri documenti sono eventualmente aggiornati.
	\item \textbf{Progettazione di dettaglio e codifica}: [Per RQualità]: vengono redatti la definizione del prodotto (approfondimento della struttura e relazione tra le varie componenti del prodotto identificate dalla specifica tecnica), codifica (sviluppo del codice in base alla definizione del prodotto), manuali (manuale per l'utente del prodotto); anche qui altri documenti vengono aggiornati e il tutto verificato.
	\item Verifica e validazione [Per RQualifica]: il prodotto viene collaudato per dimostrare che soddisfa i requisiti richiesti;
\end{itemize}

\subsection{Ciclo di vita}
Scelta del modello del ciclo di vita con motivazione. Probabilmente verrà scelto quello incrementale per cui si può prendere spunto da §2.4 di HotKeys

\subsection{Ruoli e costi}

\subsection{Scadenze}
Definizione delle scadenze delle 4 revisioni come specificate dal prof.

\begin{itemize}
	\item \textbf{22 gennaio 2016:} consegna materiale in ingresso alla \RR;
	\item \textbf{11 aprile 2016:} consegna materiale in ingresso alla  \RP;
	\item \textbf{16 maggio 2016:} consegna materiale in ingresso alla \RQ;
	\item \textbf{10 giugno 2016:} consegna materiale in ingresso alla \RA.
\end{itemize}

\newpage

\section{Analisi dei rischi}
Sezione che identifica i rischi che possono rallentare o fermare l'avanzamento del progetto. Per ogni rischio viene presentata una breve analisi, la probabilità di occorrenza e il fattore di rischio (quanto è pericoloso tale rischio).

Inoltre vengono identificati controlli (sistemi per l'analisi di questi rischi) e contromisure/prevenzioni per sistemare/prevenire tali rischi.

\subsection{Rischi del gruppo}

Rischi derivanti dai membri stessi del gruppo (mancanza di conoscenza, impegni personali, etc...).

\subsection{Rischi tecnologici}

Rischi derivati dall'utilizzo di specifiche tecnologie, librerie e strumenti. Ad esempio bug o malfunzionamenti nelle tecnologie utilizzate.

\subsection{Rischi dei requisiti}

Fraintendimenti con i requisiti o con la soluzione proposta per soddisfarli.

\subsection{Rischi organizzativi}

Rischi legati alla progettazione e organizzazione delle attività (previsioni sbagliate).

\newpage

\section{Analisi}
\subsection{Pianificazione} %dal 18/12/15 al 22/01/16
Questa attività ha inizio il 18/12/2015 e termina il 21/01/2016. I ruoli attivi in questa fase saranno:
\begin{itemize}
	\item \PM;
	\item \AM;
	\item \AN;
	\item \VR.
\end{itemize}
Le sotto attività svolte sono:
\begin{itemize}
	\item \textbf{Norme di Progetto:} attività svolta dall’\AM. Comprende la stesura di una serie di norme che dovranno essere rispettate dai componenti del gruppo per tutta la durata del progetto \prjL. Le norme sono interne al gruppo e non legate al capitolato scelto. La produzione di questo documento deve essere prioritaria per garantire uniformità e correttezza nelle successive fasi del progetto
	\item \textbf{Studio di Fattibilità:} per ogni capitolato presentato si redige una studio di fattibilità. Da questo sarà possibile capire quale progetto il gruppo deciderà di sviluppare
	\item \textbf{Analisi dei Requisiti:} gli analisti hanno il compito di ricavare tutti i requisiti necessari per la comprensione del problema e in seguito redigere il \AdR. Questo documento potrà essere aggiornato fino alla data di consegna stabilita
	\item \textbf{Piano di Progetto:} attività svolta per la pianificazione delle attività, risorse e costi della gestione del gruppo. Il \PdP{} includerà tutte queste informazioni
	\item \textbf{Piano di Qualifica:} questa attività comprende le strategie di verifica e validazione adottate. Inoltre richiede la stesura del \PdQ
	\item \textbf{Glossario:} attività svolta dai redattori parallelamente alla stesura dei documenti. Esso verrà aggiornato in modo incrementale fino al completamento della documentazione.
\end{itemize}

\begin{table}
\begin{center}
\begin{tabular}{p{0.08\textwidth} p{0.4\textwidth} p{0.2\textwidth} p{0.1\textwidth}}
\toprule
\textbf{ID}	&	\textbf{Attività}	&	\textbf{Ruolo}	&	\textbf{Ore}\\
\midrule
\midrule
\textbf{AN1} & \textbf{Norme di Progetto} &  &  \\
\midrule
AN1.1 & Stesura & Amministratore & 20 \\
\midrule
AN1.2 & Verifica & Verificatore & 7 \\
\midrule
\textbf{AN2} & \textbf{Studio di Fattibilità} & & \\
\midrule
AN2.1 & Stesura & Analista & 8 \\
\midrule
AN2.2 & Verifica & Verificatore & 3 \\
\midrule
\textbf{AN3} & \textbf{Analisi dei Requisiti} & &  \\
\midrule
AN3.1 & Stesura & Analista & 40 \\
\midrule
AN3.2 & Verifica & Verificatore & 14 \\
\midrule
\textbf{AN4} & \textbf{Piano di Progetto} & &  \\
\midrule
AN4.1 & Stesura & Project Manager \newline Amministratore & 6 \newline 10 \\
\midrule
AN4.2 & Verifica & Verificatore & 5 \\
\midrule
\textbf{AN5} & \textbf{Piano di Qualifica} & & \\
\midrule
AN5.1 & Stesura & Project Manager \newline Amministratore \newline Analista & 4 \newline 8 \newline 8 \\
\midrule
AN5.2 & Verifica & Verificatore & 7 \\
\bottomrule
\end{tabular}
\caption{Ripartizione delle risorse - Analisi}
\end{center}
\end{table}

\begin{center}
\begin{figure}
\centering
\includegraphics[width=\textwidth]{diagrammaGanttAnalisi.png}
\caption{Diagramma di Gantt della ripartizione delle risorse - Analisi}
\end{figure}
\end{center}


\newpage

\subsection{Rotazione dei Ruoli}

\begin{table}[h]
\begin{center}
\begin{tabular}{p{0.3\textwidth} p{0.05\textwidth} p{0.05\textwidth} p{0.05\textwidth} p{0.05\textwidth} p{0.05\textwidth} p{0.05\textwidth} p{0.1\textwidth}}
\toprule
\textbf{Membro}	&	\textbf{PM}	&	\textbf{Am}	& \textbf{An} & \textbf{Pg} & \textbf{Pr} & \textbf{Ve} & \textbf{Totale}\\
\midrule
\midrule
\IB & 4 & 2 & 8 &  &  & 6 & 20 \\
\midrule
\AB &  & 6 & 11 &  &  & 3 & 20 \\
\midrule
\NDC & 6 & 7 &  &  &  & 7 & 20 \\
\midrule
\TP &  & 5 & 10 & &  & 5 & 20 \\
\midrule
\WS &  & 5 & 10 & &  & 5 & 20 \\
\midrule
\AVE &  & 7 & 8 &  &  & 5 & 20 \\
\midrule
\AVI &  & 6 & 9 &  &  & 5 & 20 \\
\bottomrule
\end{tabular}
\caption{Rotazione dei ruoli - Analisi}
\end{center}
\end{table}

\begin{center}
\begin{figure}[h]
\centering
\includegraphics[width=\textwidth]{diagrammaBarreAnalisiRotazioneRuoli.png}
\caption{Diagramma della rotazione dei ruoli - Analisi}
\end{figure}
\end{center}


\newpage
\subsection{Prospetto economico}

\begin{table}[h]
\begin{center}
\begin{tabular}{p{0.2\textwidth} p{0.05\textwidth} p{0.2\textwidth} p{0.3\textwidth} }
\toprule
\textbf{Ruolo}	& \textbf{Ore} & \textbf{Costo (\euro/ora)}	& \textbf{Costo (\euro)} \\
\midrule
\midrule
Project Manager & 10 & 30 & 300\\
\midrule
Amministratore & 38 & 20 & 760\\ 
\midrule
Analista & 56 & 25 & 1.400\\ 
\midrule
Progettista & 0 & 22 & 0\\ 
\midrule
Programmatore & 0 & 15 & 0\\ 
\midrule
Verificatore & 36 & 15 & 540\\ 
\midrule
\textbf{Totale} & \textbf{140} &  & \textbf{3.000}\\
\bottomrule
\end{tabular}
\caption{Ripartizione oraria dei ruoli - Analisi}
\end{center}
\end{table}

\begin{center}
\begin{figure}[h]
\centering
\includegraphics[width=\textwidth]{diagrammaTortaAnalisiTotaleOre.png}
\caption{Diagramma della ripartizione oraria dei ruoli - Analisi}
\end{figure}
\end{center}

\newpage
\section{Progettazione architetturale} %dal 16/02/16 al 11/04/16 => 40 ore di lavoro a testa

\subsection{Pianificazione}

\subsection{Rotazione dei Ruoli}

\subsection{Prospetto economico}

\newpage
\section{Progettazione di dettaglio e codifica} % 12/04/16 al 16/05/16 => 20 ore di lavoro a testa

\subsection{Pianificazione}

\subsection{Rotazione dei Ruoli}

\subsection{Prospetto economico}

\newpage
\section{Validazione} % dal 16/05/16 al 10/06/16 => 20 ore di lavoro a testa

\subsection{Pianificazione}

\subsection{Rotazione dei Ruoli}

\subsection{Prospetto economico}


\newpage
\section{Riepilogo}


\newpage
\section{Consuntivo}

\newpage
\appendix
\section{Organigramma}

\subsection{Redazione}
Sezione che definisce chi ha redatto il documento.

\subsection{Approvazione}
Lista di chi ha approvato il documento (Project manager) + Prof

\subsection{Accettazione dei componenti}
Lista dei componenti + firme per accettare il piano.

\subsection{Componenti}
Lista dei componenti con email

\subsection{Ruoli e costi}



\end{document}
